% GENERATED FILE. Edit canonical lessons, then run npm run generate:books.
% canonical-source-hash: fnv1a64:821332eb7a08eecc
% canonical-lessons: BN-C57-gajor, BN-C57-nun, BN-C57-makhon, BN-C57-modhu, BN-C57-tel

\chapter{Carrot, Salt, Butter, Honey, Oil}
\label{ch:bn-carrot-salt-butter-honey-oil}
\hlchaptermodality{57}

\begin{chapteropening}
\textbf{By the end of this chapter:} \emph{I can say a carrot, salt, butter, honey and oil.}

Complete the last lesson of chapter 57: I can say a carrot, salt, butter, honey and oil.
\end{chapteropening}

\section[\texorpdfstring{\bn{গাজর} (gājôr)}{গাজর (gājôr)}]{\bn{গাজর} (gājôr) — a carrot}

\label{lesson:BN-C57-gajor}



\begin{quote}
Before the new one: say the Bengali for a banana, then the Bengali for a potato.
\end{quote}

\subsection*{You'll want to know: \bn{গাজর}}
\textbf{\bn{গাজর}} — \emph{gājôr} — \textquotedblleft{}a carrot\textquotedblright{}.

\begin{cousinweb}
\textbf{\bn{গাজর}} (\emph{gājôr}) is a tadbhava, a word inherited from Sanskrit and reshaped by sound change. It comes from \emph{gārjara}, \textquotedblleft{}carrot\textquotedblright{} (CDIAL 4140).
\end{cousinweb}

\begin{grammarlens}[title={What you've built}]
One more word for this chapter.
\end{grammarlens}

\subsection*{Your turn}
\begin{itemize}
\raggedright
  \item \emph{Say it:} \emph{gājôr}
  \item \emph{Say it:} \emph{gājôr}, once more
  \item \emph{Say it:} say \emph{gājôr}
  \item \emph{Recall:} say the Bengali for a banana, then the Bengali for a potato, then say \emph{gājôr} again
\end{itemize}

\begin{tcolorbox}[breakable,colback=teal!4,colframe=teal!35!black,title={Before you move on}]
What does \bn{গাজর} mean? (\textquotedblleft{}A carrot\textquotedblright{}.) Say it once more.
\end{tcolorbox}
\section[\texorpdfstring{\bn{নুন} (nun)}{নুন (nun)}]{\bn{নুন} (nun) — salt}

\label{lesson:BN-C57-nun}



\begin{quote}
Before the new one: say the Bengali for a potato, then the Bengali for a carrot.
\end{quote}

\subsection*{You'll want to know: \bn{নুন}}
\textbf{\bn{নুন}} — \emph{nun} — \textquotedblleft{}salt\textquotedblright{}.

\begin{cousinweb}
\textbf{\bn{নুন}} (\emph{nun}) goes back to Sanskrit \emph{lavaṇa}, \textquotedblleft{}salt\textquotedblright{}. CDIAL gives no Bengali form under that entry, but it lists Oriya \emph{nuṇa} there (CDIAL 10978), and Bengali \emph{nuniyā}, \textquotedblleft{}salted\textquotedblright{}, under the related \emph{lavaṇita}, \textquotedblleft{}salted\textquotedblright{} (CDIAL 10983).
\end{cousinweb}

\begin{grammarlens}[title={What you've built}]
One more word for this chapter.
\end{grammarlens}

\subsection*{Your turn}
\begin{itemize}
\raggedright
  \item \emph{Say it:} \emph{nun}
  \item \emph{Say it:} \emph{nun}, once more
  \item \emph{Say it:} say \emph{nun}
  \item \emph{Recall:} say the Bengali for a potato, then the Bengali for a carrot, then say \emph{nun} again
\end{itemize}

\begin{tcolorbox}[breakable,colback=teal!4,colframe=teal!35!black,title={Before you move on}]
What does \bn{নুন} mean? (\textquotedblleft{}Salt\textquotedblright{}.) Say it once more.
\end{tcolorbox}
\section[\texorpdfstring{\bn{মাখন} (mākhôn)}{মাখন (mākhôn)}]{\bn{মাখন} (mākhôn) — butter}

\label{lesson:BN-C57-makhon}



\begin{quote}
Before the new one: say the Bengali for a carrot, then the Bengali for salt.
\end{quote}

\subsection*{You'll want to know: \bn{মাখন}}
\textbf{\bn{মাখন}} — \emph{mākhôn} — \textquotedblleft{}butter\textquotedblright{}.

\begin{cousinweb}
\textbf{\bn{মাখন}} (\emph{mākhôn}) is a tadbhava, a word inherited from Sanskrit and reshaped by sound change. It comes from \emph{mrakṣaṇa}, \textquotedblleft{}rubbing in; ointment, oil\textquotedblright{} (CDIAL 10378).
\end{cousinweb}

\begin{grammarlens}[title={What you've built}]
One more word for this chapter.
\end{grammarlens}

\subsection*{Your turn}
\begin{itemize}
\raggedright
  \item \emph{Say it:} \emph{mākhôn}
  \item \emph{Say it:} \emph{mākhôn}, once more
  \item \emph{Say it:} say \emph{mākhôn}
  \item \emph{Recall:} say the Bengali for a carrot, then the Bengali for salt, then say \emph{mākhôn} again
\end{itemize}

\begin{tcolorbox}[breakable,colback=teal!4,colframe=teal!35!black,title={Before you move on}]
What does \bn{মাখন} mean? (\textquotedblleft{}Butter\textquotedblright{}.) Say it once more.
\end{tcolorbox}
\section[\texorpdfstring{\bn{মধু} (môdhu)}{মধু (môdhu)}]{\bn{মধু} (môdhu) — honey}

\label{lesson:BN-C57-modhu}



\begin{quote}
Before the new one: say the Bengali for salt, then the Bengali for butter.
\end{quote}

\subsection*{You'll want to know: \bn{মধু}}
\textbf{\bn{মধু}} — \emph{môdhu} — \textquotedblleft{}honey\textquotedblright{}.

\begin{cousinweb}
\textbf{\bn{মধু}} (\emph{môdhu}) is Sanskrit \emph{madhu}, \textquotedblleft{}honey, mead\textquotedblright{}, taken over whole as a learned loan, a tatsama (CDIAL 9784). The word inherited from it is \emph{mau}, which CDIAL lists.
\end{cousinweb}

\begin{grammarlens}[title={What you've built}]
One more word for this chapter.
\end{grammarlens}

\subsection*{Your turn}
\begin{itemize}
\raggedright
  \item \emph{Say it:} \emph{môdhu}
  \item \emph{Say it:} \emph{môdhu}, once more
  \item \emph{Say it:} say \emph{môdhu}
  \item \emph{Recall:} say the Bengali for salt, then the Bengali for butter, then say \emph{môdhu} again
\end{itemize}

\begin{tcolorbox}[breakable,colback=teal!4,colframe=teal!35!black,title={Before you move on}]
What does \bn{মধু} mean? (\textquotedblleft{}Honey\textquotedblright{}.) Say it once more.
\end{tcolorbox}
\section[\texorpdfstring{\bn{তেল} (tel)}{তেল (tel)}]{\bn{তেল} (tel) — oil}

\label{lesson:BN-C57-tel}



\begin{quote}
Before the new one: say the Bengali for butter, then the Bengali for honey.
\end{quote}

\subsection*{You'll want to know: \bn{তেল}}
\textbf{\bn{তেল}} — \emph{tel} — \textquotedblleft{}oil\textquotedblright{}.

\begin{cousinweb}
\textbf{\bn{তেল}} (\emph{tel}) is a tadbhava, a word inherited from Sanskrit and reshaped by sound change. It comes from \emph{taila}, \textquotedblleft{}sesame oil, oil\textquotedblright{} (CDIAL 5958). It is built on \emph{tila}, \textquotedblleft{}sesame\textquotedblright{}.
\end{cousinweb}

\begin{grammarlens}[title={What you've built}]
One more word for this chapter.
\end{grammarlens}

\subsection*{Your turn}
\begin{itemize}
\raggedright
  \item \emph{Say it:} \emph{tel}
  \item \emph{Say it:} \emph{tel}, once more
  \item \emph{Say it:} say \emph{tel}
  \item \emph{Recall:} say the Bengali for butter, then the Bengali for honey, then say \emph{tel} again
\end{itemize}

\begin{tcolorbox}[breakable,colback=teal!4,colframe=teal!35!black,title={Before you move on}]
What does \bn{তেল} mean? (\textquotedblleft{}Oil\textquotedblright{}.) Say it once more.
\end{tcolorbox}
